% This LaTeX document needs to be compiled with XeLaTeX.
\documentclass[14pt]{jsarticle}
\usepackage[top=22truemm,bottom=15truemm,left=22truemm,right=22truemm]{geometry} %ページの余白設定
\usepackage{tikz,ascmac,amsmath,amssymb,amsthm,bm,physics,arydshln,mathcomp,wrapfig,pgfplots,multicol}
\usetikzlibrary{calc}
\usetikzlibrary{intersections}
\DeclareMathOperator{\id}{id}


\begin{document}
 \ 
\bigskip

\begin{flushright}
\LARGE\textgt{第二版への序文}
\end{flushright}

\bigskip
\bigskip
\bigskip

五年前の出版の後に，本書の初版は6刷を重ねるに至った. そろそろ読者から寄せられた意見に応えて，指摘された，あるいは我々自らがみつけた誤りを 訂正し，不足を補うべきときであろう. この2版では，初版のときは時間と 分量の制限から除いた内容も追加している。これらの追加された内容は，本書 の潜在的な読者層がはつきりするにつれて，より本質的に重要であると我々が考えるようになったものである。また，我々はこの改訂でも，この教科書の本質的な性格である，厳密で，しかも自習に適切な形で素材を提供する，という 方針を保つことができたものと期待している。よって，焦点は第1版と同じく,初心者にわかりやすく，明確で，具体的な例を強調することにある。そして読者に直接，素材に関わってもらい，結果の多くを読者自身が証明するように挑戦するスタイルをとっている(解答は本書の中に与えてあるけれども!).

例と練習問題とは別にして, この版で新しく加えた内容は二つの対照的な種類のものである。その一方の新たに加えた第 7 章は，一般的な測度の間の比較 を議論するもので，内容の中心はラドンーニコデイムの定理である。ここで与 えた証明は新しいものではないが，通常の証明よりも構成的で初等的だろうと 思う。さらに, この結果を一貫した方針で用いて，直線上のルベーグースティル チエス測度の構造を考察し，符号付き測度のハーンージョルダン分解を導く. して関数と測度の変分と絶対連続性の概念が共通の源をもつことが，明確にされる. この確率論への主な応用は条件付期待値である。これについては，第 5 章で直交射影によるまた別の構成法も与えた。条件付期待値の性質は第 7 章で,離散時間マルチンゲールの初等的性質を導くために応用される.

もう一方は,各章の終わりに追加した部分である(第1章と第5章を除く). 急速に成長している数理ファナインス分野の学生が，本書の読者の大きな割合を占めるようになったことは明らかである. よって, 関連する各章の最後にこの分野からのアイデアの短い議論をおくことにした. これらの節では，本書を自己完結したものにするという目標から外れざるを得なかった. というのも, この\\
分野全体を一から展開することは,ほとんど不可能だからである. したがって,紹介するファイナンス理論の概念を，その発端から定義することも，説明することもしない. しかしその代わりに, これらの概念を測度と確率論の枠組みの中に数学的に位置づけることを目指した. これによって, ファイナンス理論の視点から書く著者が回避しがちな数学的厳密さを結論に与えることができる.

誤解を避けるため，本書の目的は測度と積分の理論のアイデアを展開するこ と，特に，確率論と(手短かにだが)ファイナンスへの応用の視点から説明する ことであることを再度，繰り返しておく。ゆえに，本書は確率論の教科書でも なければ，数理ファイナンスの教科書でもない。これらの分野の両方とも，そ れ自身の膨大な専門的文献があり，我々が本書を書いた意図は，これらの分野 を支える数学的枠組みを学生が理解するのを助けることにある。

初版における多くの間違い，誤植と概念的なものの両方を，指摘してくれた 読者達と同僚達に感謝したい。まだどうても残っている間違いについては、我々自身の責任である。速やかな訂正を容易にするため，間違いや不正確な部分を報せるためのウェブページが,

\section*{http://www.springeronline.com/1-85233-781-8}
に設けられているので，読者にはこれを参照することを勧める。また，シュプ リンガー・フェアラーク・ロンドンのステファニー・ハーディングとカレン・ ボースウイックの二人にも, この第二版を出版するにあたっての, 常なる気遣 いと手助けに感謝を述べたい。

2003 年 10 月

$$
\begin{array}{r}
\text { クラクフ(ポーランド) マレク・ツァピンスキ } \\
\text { ハル(イギリス) エクハード・コップ }
\end{array}
$$

\section*{第一版への序文}
本書で扱う中心的課題はルベーグ測度とルベーグ積分である。これらはイギ リスの学部学生向けの教程で標準的に提供される内容として, 完全に確固たる 地位を築いているとはいえない。しかし，この概念は現代の確率論を支える，抽象的測度空間の理論を展開するための主要なモデルを提供するものであり，ま た一方では，このルベーグ関数の空間は，相変わらず，関数解析およびその多 $\qquad$ フーリエ解析や偏微分方程式論 の方法論を試すための舞台の 主な供給源であり続けている.

したがって, 測度と積分の構成と性質を明確に理解しなければならないのは,解析学者の卵だけではない。数学，物理，電子工学，工学，そして最近におい てはファイナンス理論, といった幅広い数学の応用分野で, 解析的手法の応用 を研究したいと真剣に思い，基礎になっている理論を注意深く学ぶ必要のある 人々もそうなのである。

我々のみるところ, 学部学生にも学ぶことができるレベルで,このような必要を正確に満たしている教科書はほとんどない。現代確率論には多くの良い教科書があるし，そのような教科書でも，本書で展開したような道具を確実に理解することの必要性が次第に理解されてきているようである。しかし，大抵の 場合, これらの教科書で扱っている方法は学部学生には高度過ぎるか，または， どこか通り一遍でおざなりなものになっている。そのようなわけで，本書がこ のあまりにも大きい「隙間」を埋める特効薬であるとは思わないでいただきた い!!

積分の理論を展開するにあたって決めなければならない基本事項の一つは,測度と積分のどちらで始めるか，つまり，集合から始めるか，関数で始めるか である。関数解析の専門家は後者の手法を好みがちだが，確率論の展開にとっ ては明らかに前者が必要である。我々はこの議論では, 確率論研究者の側に立 つことを選び，そしてその結果として当然ながら，主な応用の場として，確率論の基本的概念と成果の系統的な展開を用いることを選んだ、テーマと用語の 現れる順序は，我々のこの選択を反映している。また，各章は確率論の概念に 関係のある，さらに進んだ内容も含んでいる。ときにこのアプローチは，その 他の方法よりも「効率」が良くないようにみえるかもしれない。しかし我々は 時々，エレガントさを犠牲にして，直接的な証明と具体的な構成をとることを 選んだ。この判断が読者の理解を増してくれるものと期待する。

紙数と時間の許す限り，積分と測度については本書だけで間に合うようにし たかったので，最終学年の学生にとっては，あまりに月並にみえる章もあるかも しれない。しかし, ハル大学で何年にも渡ってこの素材を試してきた経験によ れば,聴講生達が解析学の基本概念に十分に親しみ, それらを確かなものにし ているとは思えないことが多かった。ゆえに，本書の準備の章には，初等的実解析の基本概念の復習とともに,リーマン積分の説明なども含めることにした。

積分と測度を応用する主な分野として, 我々は確率論を選んだものの, 本書 は初等確率論の教科書ではない。そのような良い教科書は確率論分野にたくさ んある。\\
本書は進んだ解析学を扱う本ではないが，（表面だけをざっと見るのではな く) 読者にちゃんと勉強してもらうことを意図しているし，自習向けにも講義用テキストとしても便利であるように目指した。よって，結果のかなりの部分 である，「命題」とされた内容はその場で証明を与えず，その先を読む前に読者 に証明を試みてもらうため残すようにしたし(その多くは, どこから手をつけ るかのヒントをつけてある)，豊富な練習問題も含めた。自習を助けるため，命題の証明は各章の終わりに，練習問題の解答のあら筋は本書の最後に与えてあ る.よって、勤勉な読者には、ほとんど疑問は残らないだろう.

測度と積分の主な定義の動機と準備を与える，導入のための章の後，第２章 ではルベーグ測度の詳細な構成とその性質を扱い，確率空間への応用のため公理の抽象化へと進む。ここでは，後の章のための雛形を用意する．後の章では，確率論独特の内容として，独立性の概念がより一般的な文脈から追及される。

第３章は非負可測関数の積分を展開し，確率変数とそこから導かれる確率分布を導入する。一方, 第 4 章ではルベーグ積分の主な極限定理を展開し,リー マン積分と比較する。この確率論への応用として期待値を議論し, 密度や特性関数の役割をみる。

第 5 章での動機はもっと関数解析的なものである。つまり，その焦点はルベー グ関数空間で，二乗可積分関数の空間 $L^{2}$ の特別な役割の議論も含んでいる。第 6 章では測度理論に戻って, 直積測度とフビニの定理を詳細に展開する。ここ では，確率論における結合分布と条件付確率，期待値の役割に導かれる。最後 に, 可積分関数の列の収束の主要な種類の議論に続いて，第７章ではあつかま しく確率論への偏向を認め，主要な極限定理を扱って，さらに最後にはリンド バーグーフェラーの中心極限定理を示す.

本書は教科書であって論文ではないので, ここでの議論は徹底的に厳密なも のではない．本書で扱ったテーマの範囲はおそらく，第三学年の一学期で終え るには，やや量が多いだろう。しかし，最初の五つの章がこのような学生にちょ うどよい教程だろうし，最後の二つの章は自習として，または，確率論に進む ことを希望する学生のためのセミナー用として，残すことも考えられる。また は, 解析学に十分な知識のある学生向けには, 最初の二つの章を基礎知識のま とめとみなして、残りの章を一学期の教程としてもよいかもしれない。

1998 年 5 月

$$
\begin{array}{r}
\text { クラクフ(ポーランド) マレク・ツァピンスキ } \\
\text { ハル(イギリス) エクハード・コップ }
\end{array}
$$

\section*{目次}
\begin{enumerate}
  \item 動機と準備\\
1.1 記号と集合論の基礎\\
1.2 リーマン積分：適用範囲と制限 ..... 7\\
1.3 ランダムに数を選ぶとは ..... 13
  \item 測 度 ..... 15\\
2.1 零集 合 ..... 15\\
2.2 外測度 ..... 20\\
2.3 ルベーグ可測集合とルベーグ測度 ..... 25\\
2.4 ルベーグ測度の基本的性質 ..... 33\\
2.5 ボレル集合 ..... 36\\
2.6 確率 ..... 41\\
2.7 命題の証明 ..... 47
  \item 可測関数 ..... 49\\
3.1 実数直線の拡張 ..... 49\\
3.2 ルベーグ可測関数 ..... 49\\
3.3 例 ..... 52\\
3.4 性 質 ..... 53\\
3.5 確 率 ..... 58\\
3.6 命題の証明 ..... 65
  \item 積 分 ..... 67\\
4.1 積分の定義 ..... 67\\
4.2 単調収束定理 ..... 72\\
4.3 可積分関数 ..... 76\\
4.4 優収束定理 ..... 81\\
4.5 リーマン積分との関係 ..... 85\\
4.6 可測関数の近似 ..... 90\\
4.7 確率論 ..... 93\\
4.8 命題の証明 ..... 105
  \item 可積分関数の空間 ..... 109\\
$5.1 L^{1}$ 空間 ..... 110\\
$5.2 L^{2}$ ヒルベルト空間 ..... 114\\
$5.3 L^{p}$ 空 間：完備性 ..... 122\\
5.4 確率論 ..... 127\\
5.5 命題の証明 ..... 135
  \item 積測度 ..... 137\\
6.1 多次元ルベーグ測度 ..... 137\\
6.2 積 $\sigma$-加法族 ..... 138\\
6.3 積測度の構成 ..... 139\\
6.4 フビニの定理 ..... 146\\
6.5 確 率 ..... 149\\
6.6 命題の証明 ..... 160
  \item ラドンーニコディムの定理 ..... 161\\
7.1 密度と条件 ..... 161\\
7.2 ラドンーニコディムの定理 ..... 162\\
7.3 ルベーグースティルチェス測度 ..... 173\\
7.4 確 率 ..... 192\\
7.5 命題の証明 ..... 210
  \item 極限定理 ..... 216\\
8.1 収束の種類 ..... 216\\
8.2 確 率 ..... 218\\
8.3 命題の証明 ..... 251\\
解 答 ..... 253\\
補遺 ..... 267\\
参考文献 ..... 270\\
訳者あとがき ..... 271\\
索 引 ..... 273
\end{enumerate}

\section*{動機と準備}
人生には不確実性がつきまとう。計画が思いどおりにいくことは滅多にないし，その ために我々は幼い頃から，ある事柄が他に比べて「より起こりそうだ」という，出来事 が起こる確からしさについて考えるようになる。こういつた曖昧な概念を確率モデルに 変換するとは，つまり，不確実性について考えるための合理的な枠組みを構成すること である。この枠組みは一般的でなければならないし，十分な量の事前情報を吟味できる ときも，ほとんどそういった情報なしに進まなければならないときも，等しく状況を扱 えなければならない。どんな場合にもある程度の決断が必要になるが, 我々は計量的に 一般的な法則を導くことを可能にする，秩序だった理論的な枠組みと方法を探す。

こうして我々は確率の数学的なモデルへと導かれる。これは経験的観察の理想化され た抽象化で，幅広く応用可能であること，正確で単純であることの条件を満たさなけれ ばならない。本書での中心的な興味は，無限の標本空間と無限の試行列をも考えること ができるような，一般的な確率モデルを構成し，応用することである. 「[0,1] 区間からラ ンダムに一つの数字を選ふ」といった一見単純な概念を正しく意味づけしょうと試みた り，独立な試行を繰り返していった極限での振舞いを理解したいと考えたりすると,こ ういったモデルが必要になることが容易にわかる。初等的な確率が, 起こりうる出来事 の集合の大きさの比で計算されるのと全く同じように, 我々が解くべき根本的問題は，無限個の要素をもつ集合の「大きさ」を測ることになる. 少なくとも実数直線上の集合 については，基本的な実解析のアイデアが確実な答えを与えてくれるし，またこれは実 のところ，確率論の基礎に用いる抽象的な公理論的枠組に必要な全てのアイデアを含ん でもいる。この理由から，「測度」の概念，特に，実数 $\mathbb{R}$ 上の「ルベーグ測度」の概念 を展開することが，本書の中心部分になる。

\section*{1.1 記号と集合論の基礎}
測度論の典型的な場合では, 任意に与えられた部分集合の族を扱い, この族 に属する集合に対して実数を割り当てるような関数を考える。よって, 基本的 な集合論の記号や集合の操作について，振り返っておく必要がある. 同様に,可算無限集合と非可算無限集合，特に実数直線 $\mathbb{R}$ 上の部分集合を扱うが，それら

の違いも議論する。また他に, 数列, 級数, 開集合の極限のような解析学の概念も必要になる。読者はこれらをよく知っているものと仮定しているので, 本章については軽く読み流す程度でもよい。実際本章は，使用する記号について 説明し，本書を自己完結させるためのものであり，自習に便利だろう.ここで の議論はかなりくだけたもので，根本的な問題には言及していない。よって読者は、その証明等詳細を知りたいときには，解析学の基本的な教科書を参照す るように. ここでは, 入門的教科書 [8] と [11] を挙げるだけにしておく。

\subsection*{1.1.1 集合と関数}
本書では常に, 集合の操作をある一般的な集合 $\Omega$ の部分集合の操作として扱 う. この集合 $\Omega$ 自体の素性は文脈から明らかになるが, 多くの場合, 実数全体 $\mathbb{R}$ か，その部分集合である。我々は「集合」の概念を未定義のままにし，既に 与えられたものと考え，集合の要素と操作についてだけ考えることにする．要素を全くもたない集合，空集合を と書く. 集合は通常, $A, B$ といった大文字 で書くことにする。

ある集合の要素(元)であることは, $\in$ で表す. つまり, $x \in A$ は, $x$ が集合 $A$ の要素であることを意味する。集合の包含関係 $A \subset B$ は, $A$ の全ての要素が $B$ の要素であることを意味する。これは $A$. $B$ が等しい場合も含んでいる. 包含関係が強い意味である場合，つまり， $A \subset B$ であるが， $B$ に含まれるが $A$ には 含まれない要素 $(x \notin A$ と書く)が存在する場合には, そのように明記する. 性質 $P$ をもつような $A$ の要素の集合を, 記号 $\{x \in A: P(x)\}$ で表す. $A$ の全て の部分集合の集合（つまり，Aのべき集合)を, $\mathcal{P}(A)$ と書く.

集合の積集合 (共通集合)を $A \cap B=\{x: x \in A$ かつ $x \in B\}$ で定義し, 和集合を $A \cup B=\{x: x \in A$ または $x \in B\}$ で定義する. 集合 $A$ の補集合 $A^{c}$ とは, $A$ の要素でない $\Omega$ の要素で構成される集合のことである. 同じことを, $A^{c}=\Omega \backslash A$ とも書く.より一般に, 集合の差 (差集合) $B \backslash A=\{x \in B: x \notin A\}=B \cap A^{c}$

と, 対称差 $A \triangle B=(A \backslash B) \cup(B \backslash A)$ も定義する. $A=B$ のとき，かつそのとき に限り, $A \triangle B=\emptyset$ であることに注意せよ.

積集合 (resp. 和集合)は論理連結語「かつ」(resp.「または」) で表現されて いて，論理記号 (存在する) と $\forall$ (全ての)を用いて，任意の族に拡張すること ができる。つまり，ある集合 1 を添字集合として,

$$
\bigcap_{\alpha \in \Lambda} A_{\alpha}=\left\{x: \text { 全ての } \alpha \in \Lambda \text { に対し } x \in A_{\alpha}\right\}=\left\{x:{ }^{\forall} \alpha \in \Lambda, x \in A_{\alpha}\right\} \text {, }
$$

$$
\bigcup_{\alpha \in \Lambda} A_{\alpha}=\left\{x: \text { ある } \alpha \in \Lambda \text { に対し } x \in A_{\alpha}\right\}=\left\{x:{ }^{\exists} \alpha \in \Lambda, x \in A_{\alpha}\right\}
$$

と定義する. これらはド・モルガン (de Morgan) の法則で以下のように結びつ いている：

$$
\left(\bigcup_{\alpha} A_{\alpha}\right)^{c}=\bigcap_{\alpha} A_{\alpha}{ }^{c}, \quad\left(\bigcap_{\alpha} A_{\alpha}\right)^{c}=\bigcup_{\alpha} A_{\alpha}{ }^{c}
$$

もし, $A \cap B=\emptyset$ ならば， $A$ と $B$ は非交叉的である，あるいは排他的である，互 いに素である，等という。集合の族 $\left(A_{\alpha}\right)_{\alpha \in \Lambda}$ は, 全ての $\alpha \neq \beta(\alpha, \beta \in \Lambda)$ に対し $A_{\alpha} \cap A_{\beta}=\emptyset$ であるとき, 互いに非交叉的, 互いに排他的，あるいは互いに素 であるという.

集合 $A$ と $B$ の直積 $A \times B$ とは, 順序を考慮した対の集合 $A \times B=\{(a, b)$ : $a \in A, b \in B\}$ のことである。既に用いたように, 扱う数の基本系として自然数,整数，有理数，実数をそれぞれ $\mathbb{N}, \mathbb{Z}, \mathbb{Q}, \mathbb{R}$ の記号で表す．実数 $\mathbb{R}$ の区間は その両端点と, 端点を含む場合には鍵括弧 $([,]$) で, 含まない場合には丸括弧$ ((,)$) で表す. 例えば, $[a, b)=\{x \in \mathbb{R}: a \leq x<b\}$ 等とする。また，記号 $\infty$ と $-\infty$ を非有界な区間を表すのに用いる。例えば, $(-\infty, b)=\{x \in \mathbb{R}: x<b\}$ や， $[0, \infty)=\{x \in \mathbb{R}: x \geq 0\}=\mathbb{R}^{+}$等とする。集合 $\mathbb{R}^{2}=\mathbb{R} \times \mathbb{R}$ で平面を, さら に一般的には， $\mathbb{R}^{n}$ で $\mathbb{R}$ の自身との $n$ 回直積を表す. つまり， $n$ 個の実数の組 $\left(x_{1}, \ldots, x_{n}\right)$ の全体の集合のことである。区間の直積も同様に記述し，長方形と よぶ.

形式的には, 関数 $f: A \rightarrow B$ とは, 一つ目の要素が二つ目を要素を決定し ている，つまり, $(a, b),(a, c) \in f$ ならば $b=c$ であるような, $A \times B$ の部分集合のことである。その定義域 $\mathcal{D}_{f}=\left\{a \in A:{ }^{\exists} b \in B,(a, b) \in f\right\}$ と値域 $\mathcal{R}_{f}=\left\{b \in B:{ }^{\exists} a \in A,(a, b) \in f\right\}$ が, その関数の働く範囲を表している。形式ばらずにいえば，関数 $f$ は $A$ の各要素に対して，Bの要素を対応させるも ので, 各 $a \in A$ は高々一つの像 $b \in B$ をもつ。 よって, これを $b=f(a)$ と書 く. 集合 $X \subset A$ はその像 $f(X)=\{b \in B:$ ある $a \in X$ に対して $b=f(a)\}$ をも ち, 集合 $Y \subset B$ の逆像とは $f^{-1}(Y)=\{a \in A: f(a) \in Y\}$ のことである。関数 $f_{1}: A \rightarrow B$ と $f_{2}: B \rightarrow C$ の合成 $f_{2} \circ f_{1}$ とは, $h(a)=f_{2}\left(f_{1}(a)\right)$ で定義される関数 $h: A \rightarrow C$ のことである. $A=B=C$ のときには, $x \mapsto\left(f_{1} \circ f_{2}\right)(x)=f_{1}\left(f_{2}(x)\right)$ も, $x \mapsto\left(f_{2} \circ f_{1}\right)(x)=f_{2}\left(f_{1}(x)\right)$ も, $A$ から $A$ への関数を定義する. 一般には, この両者は同じでない。例えば, $f_{1}(x)=\sin x, f_{2}(x)=x^{2}$ とすると, $x \mapsto \sin \left(x^{2}\right)$ と $x \mapsto(\sin x)^{2}$ は等しくない.\\
関数 $g$ は, もし $\mathcal{D}_{f} \subset \mathcal{D}_{g}$ かつ $\mathcal{D}_{f}$ 上で $g=f$ ならば, $f$ の拡張または延長で あるという。逆に, $f$ は $g$ の $\mathcal{D}_{f}$ への制限であるという. これらの概念は実数 に値をとる集合関数に対し頻繁に用いられ，その場合，考える定義域は集合の 族，値域は $\mathbb{R}$ の部分集合である。

実関数の代数は各点ごとの意味で定義する。つまり, 関数 $f$ と関数 $g$ の和 $f+g$ と積 $f g$ はそれぞれ,

$$
(f+g)(x)=f(x)+g(x), \quad(f g)(x)=f(x) g(x)
$$

で与えられる。

集合 $A$ の定義関数 $\mathbb{1}_{A}$ とは関数

$$
\mathbb{1}_{A}(x)= \begin{cases}1 & (x \in A \text { の場合 }) \\ 0 & (x \notin A \text { の場合 })\end{cases}
$$

のことである。以下の性質,

$$
\mathbb{1}_{A \cap B}=\mathbb{1}_{A} \mathbb{1}_{B}, \quad \mathbb{1}_{A \cup B}=\mathbb{1}_{A}+\mathbb{1}_{B}-\mathbb{1}_{A} \mathbb{1}_{B}, \quad \mathbb{1}_{A^{c}}=1-\mathbb{1}_{A}
$$

に注意せよ。

さらにもう一つ基本的な集合論の概念を必要とするが，既に馴染み深いもの だろう。任意の集合 $E$ にいて, $E$ 上の同値関係とは，以下の三つの性質を満 たす関係（つまり，E×Eの部分集合 $R$ で, $(x, y) \in R$ を $x y$ と書く)のことで ある：

\begin{enumerate}
  \item 反射的：全ての $x \in E$ に対し， $x \sim x$,

  \item 対称的： $x \sim y$ ならば $y \sim x$,

  \item 推移的： $x \sim y$ かつ $y \sim z$ ならば, $x \sim z$.

\end{enumerate}

集合 $E$ 上の同値関係は $E$ を非交叉的な同値類に分割する。任意の $x \in E$ に 対し, $[x]=\{z \in E: z \sim x\}$ と書き, これを $x$ の同値類という. これはつま り, $x$ と同値な $E$ の要素全ての集合のことである. よって $x \in[x]$ であり, ゆ えに, $E=\bigcup_{x \in E}[x]$ である. これは非交叉的な集合和になっている。実際, もし $[x] \cap[y] \neq \emptyset$ ならば, $x \sim z$ かつ $z \sim y$ であるような $z \in E$ が存在するので, $[x]=[y]$ である. 全ての同値類の集合を $E /$ 〜 と書く.

\subsection*{1.1.2 $\mathbb{R}$ の可算集合と非可算集合}
集合 $A$ と $\mathbb{N}$ の部分集合との間に 1 対 1 対応があるとき, つまり,異なる点 を異なる点に写す関数 $f: A \rightarrow \mathbb{N}$ があるとき, 集合 $A$ は可算であるという。形\\
式的な議論は避けるが，この対応が $\mathbb{N}$ の最初の部分 $\{1,2, \ldots, N\}$ だけで作られ るとき $(N \in \mathbb{N}), A$ は有限であるといい, $\mathbb{N}$ 全体が必要なときは, $A$ は可算無限であるという。可算集合の可算個の和が再び可算であることはすぐにわかる. よって, 特に有理数の集合 $\mathbb{Q}$ も可算である.

カントールは, 集合 $\mathbb{R}$ は $\mathbb{N}$ (の部分集合) との間に 1 対 1 対応を作ることが できないことを示した。これが，非可算集合の一例である。カントールの証明 では，各実数が一意に十進法で表現できるとする。(実際，常に項が無限に続く 表現のほうを選べばよい。）また，区間 $[0,1]$ が非可算であることを示すことに 制限してもよい。（これはなぜか？)

もし, この集合が可算であったならば，その要素を列 $\left(x_{n}\right)_{n \geq 1}$ として書くこと ができる。そして，各 $x_{n}$ は一意な十進法表現

$$
x_{n}=0 . a_{n 1} a_{n 2} a_{n 3} \ldots a_{n n} \ldots
$$

をもつ. ここで，各 $a_{i j}$ は集合 $\{0,1,2, \ldots, 9\}$ の要素である. よって, これらを 以下のように並べることができる：

$$
\begin{aligned}
& x_{1}=0 . a_{11} a_{12} a_{13} \ldots \\
& x_{2}=0 . a_{21} a_{22} a_{23} \ldots \\
& x_{3}=0 . a_{31} a_{32} a_{33} \ldots
\end{aligned}
$$

いま, $a_{n n}$ と異なる数字 $b_{n}$ を用いて, $y=0 . b_{1} b_{2} b_{3} \ldots$ とおこう. この十進法表現はある数 $y \in[0,1]$ を定義するが, どの $x_{n}$ とも異なる。（なぜなら， $x_{n}$ とは $n$ 桁目の数字 $x_{n n}$ が異なるから。) ゆえに，最初に並べられた列は $[0,1]$ を尽く すことはできず，この矛盾から [0,1] は可算ではありえない.

二つの可算集合の和は可算であり, そして $\mathbb{Q}$ も可算であるから, $\mathbb{R} \backslash \mathbb{Q}$ は非可算である。つまり，有理数よりも無理数のほうがずっと「多い」！これをもっ と理解しやすくする方法の一つは， $\mathbb{R}$ の区間からランダムに数を選ぶ問題を考 えることだろう.

有理数は，その十進法表現が循環する(ある桁から先，決まった有限の表現の 繰返しになる) 実数であることを思い出そう。（十進法表現が有限で「終了」す る場合も「循環」に含める.) ここで, [0,1] 区間からランダムに実数を一つ選 ぶ，という操作を想像してみよう． $\mathbb{R}$ を全ての実数を含む池だと考えて, この 池から一回に一匹ずつ数を「釣り上げる」と想像すればよい。

最初の数字が有理数である，つまりその十進法表現が循環している数である\\
可能性はどれくらいだろうか？これは十面のサイコロを無限回投げて，ある有限回の後は，全ての部分列で同じ数字が出る，と仮定するのと同じことである. これは全く起こりそうにない。ゆえに,ここで用いたような「ランダムな」,「偏 りのない」または，一様な方法で実数直線上の集合の大きさを測るときには，可算集合 $(\mathbb{Q}$ を含め)を「無視する」ことができる，ということに，それほど驚 くことはないだろう。しかしながら，おそらく，ここでの視点からすれば驚く べきことはむしろ，後にみるように，非可算の集合なのに「無視」できるもの がある, という事実のほうだろう.

\subsection*{1.1.3 $\mathbb{R}$ の集合の位相的性質}
開集合 $O \subset \mathbb{R}$ の定義を思い出そう.

定義 1.1. 実数直線 $\mathbb{R}$ の部分集合 $O$ が開区間の和であるとき，すなわち，あ る添字集合 $\Lambda$ (可算でも非可算でもよい)を用いて

$$
O=\bigcup_{\alpha \in \Lambda} I_{\alpha}
$$

と表されるとき，O は開集合であるという。開集合の補集合を閉集合という。 $\mathbb{R}^{n}(n>1)$ の開集合は，開区間の $n$ 重直積の和として定義する.

$\mathbb{R}$ 上では可算和を考えれば常に十分なので, この定義はその実際よりも必要以上に一般的にみえる。しかし，一般的な和を使えることが，後の議論で便利 になる。 $\Lambda$ を添字の集合として，各 $\alpha \in \Lambda$ に対し $I_{\alpha}$ が開区間であるとすると， その和が $\bigcup_{\alpha \in \Lambda} I_{\alpha}$ に等しいような可算な族 $\left(I_{\alpha_{k}}\right)_{k \geq 1}$ が存在する。しかも，区間の 列はそれぞれ互いに排他的であるように選べる.

開集合の有限個の積集合がまた開集合であることは簡単にわかる。しかし ながら，開集合の可算個の積は必ずしも開集合でない. 実際, $n \geq 1$ に対し, $O_{n}=(-1 / n, 1)$ とおけば, $E=\bigcap_{n=1}^{\infty} O_{n}=[0,1)$ となって, 開集合でない.

$\mathbb{R}$ は, $\mathbb{R}^{n}$ やり一般的な空間と違って, 線型順序をもつ. つまり, 任意に与 えられた $x, y \in \mathbb{R}$ に対して， $x \leq y$ であるか $y \leq x$ を決定できる。よって，集合 $A \subset \mathbb{R}$ について, 全ての $a \in A$ に対し $a \leq u$ とるものとして, 上界 $u$ を定義できる。下界も同様に定義する。全ての上界の最小値を上限（または，最小上界) といい, $\sup A$ と書く。下限(または, 最大下界) $\inf A$ は, 全ての下界の 最大値で定義する。 $\mathbb{R}$ の完備性とは，上から有界な全ての集合が上限をもつと いう主張のことである.\\
実数値関数 $f$ は, 各開集合 $O$ に対し $f^{-1}(O)$ も開集合であるとき, 連続であ るという．有界な閉集合上で定義された全ての連続な実数値関数は、この集合上で最大値と最小値を与える。例えば, $f:[a, b] \rightarrow \mathbb{R}$ が連続ならば, ある点 $x_{\max }, x_{\min } \in[a, b]$ で,

$$
\begin{aligned}
& M=\sup \{f(x): x \in[a, b]\}=f\left(x_{\max }\right) \\
& m=\inf \{f(x): x \in[a, b]\}=f\left(x_{\min }\right)
\end{aligned}
$$

となる。中間値の定理によれば，連続関数はこの最大最小の両端の間の任意の 中間の値をとる. つまり, 各 $y \in[m, M]$ に対し，ある $\theta \in[a, b]$ が存在して $y=f(\theta)$ となる.

実数列 $\left(x_{n}\right)$ を特徴づけるため，上極限 $\lim \sup x_{n}$ を

$$
\inf \left\{\sup _{m \geq n} x_{m}: n \in \mathbb{N}\right\}
$$

で定義し，下極限 $\liminf _{n} x_{n}$ を

$$
\sup \left\{\inf _{m \geq n} x_{m}: n \in \mathbb{N}\right\}
$$

で定義する。

数列 $x_{n}$ は, この二つの量が一致するとき, そしてそのときに限り収束し, そ の共通の値がその収束の極限である。数列 $\left(x_{n}\right)$ が収束することと, 実数 $x$ がそ の極限であることを振り返っておこう。全ての $\varepsilon>0$ に対し，ある $N \in \mathbb{N}$ が 存在して, $n \geq N$ ならば常に $\left|x_{n}-x\right|<\varepsilon$ となるとき, $x$ は $\left(x_{n}\right)$ の極限であり, $x=\lim _{n \rightarrow \infty} x_{n}$ と書くのだった. 級数 $\sum_{n \geq 1} a_{n}$ は, その部分和の列 $x_{m}=\sum_{n=1}^{m} a_{n}$ が収束 するとき収束し，その極限が級数の和 $\sum_{n=1}^{\infty} a_{n}$ である。

\section*{1.2 リーマン積分：適用範囲と制限}
本節では、入門的な解析学の教程の主要な部分をなす, リーマン積分の概観 を手短かに与える。そして、より進んだ応用のためには、,どうしてそれだけで は十分ではないのか, いくつかの理由も考えよう.

$a, b$ を $a<b$ であるような実数として, $f:[a, b] \rightarrow \mathbb{R}$ を有界な実数値関数と する. 区間 $[a, b]$ の分割とは, 有限集合 $P=\left\{a_{0}, a_{1}, a_{2}, \ldots, a_{n}\right\}$ であって,

$$
a=a_{0}<a_{1}<a_{1}<\cdots<a_{n}=b
$$

となっているもののことである。分割 $P$ に対し, 上リーマン和と下リーマン 和がそれぞれ,

$$
U(P, f)=\sum_{i=1}^{n} M_{i} \Delta a_{i}, \quad L(P, f)=\sum_{i=1}^{n} m_{i} \Delta a_{i}
$$

と定まる。ここに, $\Delta a_{i}=a_{i}-a_{i-1}$ で，各 $i \leq n$ に対して,

$$
M_{i}=\sup _{a_{i-1} \leq x \leq a_{i}} f(x), \quad m_{i}=\inf _{a_{i-1} \leq x \leq a_{i}} f(x)
$$

である.（f は各区間 $\left[a_{i-1}, a_{i}\right]$ 上で有界であるから, $M_{i}$ と $m_{i}$ は正しく定義され た実数であることに注意せよ.)

関数 $f$ のリーマン積分を定義するには，まず，任意に与えられた分割 $P$ に対 して, $L(P, f) \leq U(P, f)$ であることを示す. そして次に，任意の細分，つまり，分割 $P^{\prime} \supset P$ に対し, $L(P, f) \leq L\left(P^{\prime}, f\right)$ かつ $U\left(P^{\prime}, f\right) \leq U(P, f)$ を示さなければ ならない。最後に, 任意の二つの分割 $P_{1}$ と $P_{2}$ に対して，その和 $P_{1} \cup P_{2}$ が両者の細分であることから，任意の分割 $P, Q$ に対し， $L(P, f) \leq U(Q, f)$ であるこ とがわかる. これより，集合 $\{L(P, f) : P \text{は区間} [a, b] \text{の分割} \}$ は $\mathbb{R}$ 上で上から 有界であり，その上限を区間 $[a, b]$ 上の $f$ の下積分とよび， $\int_{a}^{b} f$ と書く. 同様 にして，上リーマン和の集合の下限を上積分とよび， $\overline{\int_{a}^{b}} f$ と書く。そして, こ の両者が一致するとき，関数 $f$ は区間 $[a, b]$ 上でリーマン可積分であるという。 その共通の値を $f$ のリーマン (Riemann) 積分とよび， $\int_{a}^{b} f$ と書くが，よく見 られる別の書き方は

$$
\int_{a}^{b} f(x) d x
$$

だろう。この定義は，特定の関数の可積分性を調べるのには便利とはいえない. しかし，以下に述べる枠組みは可積分性の便利な判定基準になる。証明につい ては，例えば[8]を参照せよ.

定理 1.2 (リーマンの判定条件). 関数 $f:[a, b] \rightarrow \mathbb{R}$ は, 全ての $\varepsilon>0$ に対し, ある分割 $P_{\varepsilon}$ が存在して, $U\left(P_{\varepsilon}, f\right)-L\left(P_{\varepsilon}, f\right)<\varepsilon$ であるとき，かつそのときに 限り,リーマン可積分である.

例 1.3. $f(x)=\sqrt{x}$ のときに $\int_{0}^{1} f(x) \mathrm{d} x$ を計算してみよう. 直ちに問題にな るのは，平方数以外ではその平方根をみつけるのが難しいことである。よって，異なる区間の区間長さでは同じことがいえないことになるが，分割点を平方数 だけに限ることにする。(しばしばこういう制限で計算が簡単になるにしても,\\
全ての分割で成立していなければならないことはいうまでもない.）実際，分割を

$$
P_{n}=\left\{0,\left(\frac{1}{n}\right)^{2},\left(\frac{2}{n}\right)^{2}, \ldots,\left(\frac{i}{n}\right)^{2}, \ldots, 1\right\}
$$

のようにとって，上リーマン和と下リーマン和を考える。 $f$ が増加関数である ことより，

$$
\begin{aligned}
& U\left(P_{n}, f\right)=\sum_{i=1}^{n}\left(\frac{i}{n}\right)\left\{\left(\frac{i}{n}\right)^{2}-\left(\frac{i-1}{n}\right)^{2}\right\}=\frac{1}{n^{3}} \sum_{i=1}^{n}\left(2 i^{2}-i\right) \\
& L\left(P_{n}, f\right)=\sum_{i=1}^{n}\left(\frac{i-1}{n}\right)\left\{\left(\frac{i}{n}\right)^{2}-\left(\frac{i-1}{n}\right)^{2}\right\}=\frac{1}{n^{3}} \sum_{i=1}^{n}\left(2 i^{2}-3 i+1\right)
\end{aligned}
$$

となるから，

$$
U\left(P_{n}, f\right)-L\left(P_{n}, f\right)=\frac{1}{n^{3}} \sum_{i=1}^{n}(2 i-1)=\frac{1}{n^{3}}\{n(n+1)-n\}=\frac{1}{n}
$$

よってこの差は，nを十分大きく選ぶことで，与えられた任意の $\varepsilon>0$ より小 さくできる。そして, $U\left(P_{n}, f\right)$ と $L\left(P_{n}, f\right)$ の両方が同じく, $2 / 3$ に収束するこ とも容易にわかるから，求める積分は $2 / 3$ である。

リーマンの判定条件も, リーマン可積分であるクラスの正確な描像を与え てくれるものではない。しかしながら，任意の有界な単調関数がこのクラスに 属することは容易にわかるし( [8] 参照), もう少し難しいが, 任意の連続関数 $f:[a, b] \rightarrow \mathbb{R}$ (よって自動的に有界関数)が，リーマン可積分であることも示す ことができる。

これは多くの実際的な目的においては十二分な情報であるし，上の例で示し たような面倒な計算は以下の定理で避けることができる.

定理 1.4 (微分積分学の基本定理). 関数 $f:[a, b] \rightarrow \mathbb{R}$ が連続で, 関数 $F:$ $[a, b] \rightarrow \mathbb{R}$ の微分が $f$ となっている（つまり， $(a, b)$ 上で $\left.F^{\prime}=f\right)$ ならば,

$$
F(b)-F(a)=\int_{a}^{b} f(x) d x
$$

が成立する。

この結果によって, リーマン積分は微分と結びつき, Fがf の原始関数(ま たは、「逆微分」) であること，つまり定数の差を許して，

$$
F(x)=\int_{a}^{x} f(u) d u
$$

であることが示され, よって, 解析学の教程で教わる積分計算の技術が正当化 される.

この連続性の条件は緩和することができる。最初にすることは，容易にわか るように, $f$ は有限個の点を除いて $[a, b]$ 上で有界かつ連続であると仮定する ことである。このとき，fはリーマン可積分になる。区間を分割し，分割され たそれぞれの上では $f$ が連続であるような場合を考えれば、これが成立するこ とがわかる. 各区間上では $f$ は可積分だから，もとの区間全体でも可積分であ る. 一つの例として, $[0,1]$ 上の関数 $f$ で，点 $a_{1}, \ldots, a_{n}$ では 1 の値をとり，そ れ以外の全ての点で 0 に等しいようなものを考えよう.これは $[0,1]$ 上で可積分で，その積分は０である。

しかしながら、これをさらに押し進めていくには、ルベーグ積分の力が必要 になる。定理 4.33 でみるように, $f$ がリーマン可積分であることと， $f$ が $[a, b]$上の「殆ど(ほとんど)至る所」で連続であることとは同値である。この結果は 全く自明でない。例えば,ディレクレ(Dirichlet)による，以下のような関数 $f$ が $[0,1]$ 上でリーマン可積分であることを直接に示そうと試みてみれば，その 困難さに気づくだろう：

$$
f(x)= \begin{cases}\frac{1}{n} & \left(x=\frac{m}{n} \in \mathbb{Q} \text { の場合 }\right) \\ 0 & (x \notin \mathbb{Q} \text { の場合 })\end{cases}
$$

事実，この関数 $f$ が各無理数点で連続であること，全ての有理数点で非連続で あることを示すことは、難しくない([8] 参照). えに(後でわかるように), こ の関数は区間 $[0,1]$ 上の「殆ど至る所」で連続なのである。

本書の目的は積分の「ルベーグの理論」を説明することなので,いったいどう

して, この新たな理論が必要になるのか, 議論しておく必要があるだろう.上 で示した単なるリーマン積分では，何か都合の悪いことがあるのだろうか？

第一に，その適用範囲：リーマン積分は我々が扱いたい全ての種類の積分 を扱えない。

すぐにわかるいくつかの結果は，有界区間上の連続関数の性質に依るもので ある。有界でない区間上の積分，例えば，

$$
\int_{-\infty}^{\infty} \mathrm{e}^{-x^{2}} \mathrm{~d} x
$$

や, 非有界な関数の積分

$$
\int_{0}^{1} \frac{1}{\sqrt{x}} d x
$$

等を扱うためには，極限の議論で定義される、「特異」リーマン積分に頼らなけ ればならない。例えば，上の例では，

$$
\int_{-n}^{n} \mathrm{e}^{-x^{2}} \mathrm{~d} x, \quad \int_{\varepsilon}^{1} \frac{1}{\sqrt{x}} \mathrm{~d} x
$$

を考えて, $n \rightarrow \infty$ や $\varepsilon \rightarrow 0$ とするのである。 これは，必ずしも深刻な欠点で はない。

第二に, 区間への依存：つまり，区間とは異なる一般的な集合上での積分 や，関数の値がそのような集合上に「やっかいに」分布する場合の積分を扱う， やさしい方法がない。例えば，区間 $[0,1]$ 上の有理数 $\mathbb{Q}$ の定義関数 $\mathbf{1}_{\mathbb{Q}}$ の上リー マン和と下リーマン和を考えてみよう。どう $[0,1]$ を分割してみても，各分割区間は有理数点と無理数点の両方を含むので，各上リーマン和は 1 で，下リー マン和は０である。ゆえに，区間 $[0,1]$ 上で $f$ のリーマン積分は計算できない. この関数は単純に，「不連続すぎる」のだ。（読者は，この $f$ が $[0,1]$ 上の全ての 点で不連続であることを，容易に納得できるだろう。）

第三に, 完備性の欠如：応用の視点から考えてさらに深刻な問題は,リー マン積分が関数の列の極限をとる操作と相性が悪いことである。次のような命題「リーマン可積分な関数の列 $\left(f_{n}\right)$ が (ある適切な意味で) $f$ に収束するなら

ば， $\int_{a}^{b} f_{n} \mathrm{~d} x \rightarrow \int_{a}^{b} f \mathrm{~d} x$ である」が成立するのを期待するのは自然なことであ ろう.これが成立しないような二つの例を挙げ，関数の列 $\left(f_{n}\right)$ が各点で $f$ に収束する，つまり，全ての $x$ について $f_{n}(x) \rightarrow f(x)$ であるときにどのような困難 が生じうるかを示す.

\begin{enumerate}
  \item この極限は必ずしもリーマン可積分でない。ゆえに，収束性の疑問は意味さえなさない。例えば, $f=1_{\mathbb{Q}}, f_{n}=1_{A_{n}}$ ととればよい。ここで, $A_{n}=$ $\left\{q_{1}, \ldots, q_{n}\right\}$ で，数列 $\left(q_{n}\right)(n \geq 1)$ は有理数を一列に並べたものである。こ の場合, $\left(f_{n}\right)$ 列は単調増加でさえある。

  \item 極限はリーマン可積分であるが，リーマン積分の収束は成立しない例. $f=0$ として, 区間 $[a, b]$ は $[0,1]$ としよう. そして,

\end{enumerate}

$$
f_{n}(x)=\left\{\begin{array}{cl}
4 n^{2} x & \left(0 \leq x<\frac{1}{2 n} \text { の場合 }\right) \\
4 n-4 n^{2} x & \left(\frac{1}{2 n} \leq x<\frac{1}{n} \text { の場合 }\right) \\
0 & \left(\frac{1}{n} \leq x \leq 1 \text { の場合 }\right)
\end{array}\right.
$$

とおく。これはその積分が常に1の連続関数である。一方で，全ての $x$ に


図 $1.1 f_{n}$ のグラフ

ついて十分大きい $n$ に対し $f_{n}(x)=0$ だから $(1 / n<x$ であるようとればよ い), 列 $f_{n}(x)$ は $f=0$ に収束する。（図1.1を見よ.)

一様収束の概念を導入することによってこの種の問題を避けることができる. すなわち, $C([0,1])$ の列 $\left(f_{n}\right)$ が $f$ に一様に収束するとは，数列

$$
a_{n}=\sup \left\{\left|f_{n}(x)-f(x)\right|: 0 \leq x \leq 1\right\}
$$

が 0 に収束することである，とするのである。このときにはリーマン積分が収束すること，つまり,

$$
\int_{0}^{1} f_{n}(x) \mathrm{d} x \rightarrow \int_{0}^{1} f(x) \mathrm{d} x
$$

を容易に証明することができる。

しかしながら, 「距離」 $\sup \{|f(x)-g(x)|: 0 \leq x \leq 1\}$ はこのような積分自体 とは何の関係もないし，一様収束は多くの場合，あまりに強い制限である。よ り自然な「距離」の概念の一つは $\int_{0}^{1}|f(x)-g(x)| \mathrm{d} x$ であるが、これもまた別 の問題を生じる。実際,

$$
g_{n}(x)=\left\{\begin{array}{cl}
0 & \left(0 \leq x \leq \frac{1}{2} \text { の場合 }\right) \\
n\left(x-\frac{1}{2}\right) & \left(\frac{1}{2}<x<\frac{1}{2}+\frac{1}{n} \text { の場合 }\right) \\
1 & (\text { その他 })
\end{array}\right.
$$

としてみると, $m, n \rightarrow \infty$ のとき, $\int_{0}^{1}\left|g_{n}(x)-g_{m}(x)\right| \mathrm{d} x \rightarrow 0$ が容易に示せる. つまり，図 1.2 で影をつけた部分が消える。（このとき $\left(f_{n}\right)$ はこの距離の下で コーシー(Cauchy) 列であるという.) しかしながら，各点の極限は $x>1 / 2$ の とき $f(x)=1$ で，他では 0 だから，この列の収束先は $f=\mathbb{1}_{\left(\frac{1}{2}, 1\right]}$ となり，これ は連続関数ではない。よって, 全ての連続関数 $f:[0,1] \rightarrow \mathbb{R}$ の空間 $C([0,1])$ は, この観点からは小さすぎることになる。


図 $1.2 \quad g_{n}, g_{m}$ のグラフ

この状況は, 有理数の全体だけからなる集合 $\mathbb{Q}$ だけでは足りず, $\mathbb{R}$ を考え なくてはならなかったこと（コーシー列が $\mathbb{Q}$ の中に極限をもたないことがある.例えば， $\sqrt{2}$ の有理数近似の列等)に似ている。実数 $\mathbb{R}$ の場合の完備性の重要ざ を思い出せば，このような欠如のない積分の理論を探すことは自然だろう.こ の過程で, 我々は新しい理論, つまり, リーマン積分をその特別な場合として 含み, 上に挙げたような問題を解決する理論に出会うことになる。

\section*{1.3 ランダムに数を選ぶとは}
ルベーグ測度の理論を展開し，それによって実数 $\mathbb{R}$ のあらゆる部分集合につ いて，それらの「長さ」を意味づけする前に少し休想して，我々がこの理論を 求めるいくつかの実際的な動機について考えておこう. 有限の標本空間での初等的な確率の単純さは，例えば「0と1 の間の数をランダムに選ぶ」といった,結果に無限の可能性がある場合を考えた途端に消え失せてしまう。与えられた $x \in[0,1]$ が選ばれる「確率」の意味を正当化しなければならないからである.同様の，しかしもう少し一般的な質問は以下のようなものである.「選んだ数が 有理数である確率はいくらか？」

最初の基本的な疑問は，ランダムに数 $x$ を選ぶ，というとき，それが何を意味しているのか, ということである。「ランダムに」という言葉は, このような 各試行について, 各実数が「同様に確からしく」選ばれること，つまり， [0,1]上の一様確率分布を定義することを意味しているように思われる。しかし, 選 ばれる「数」の可能性は無限にある。ゆえに, 与えられた $x$ とつ固定したと きその $x$ が選ばれるという事象 $A_{x}$ の確率は０でなければならない.しかしなが ら一方では，０と1の間のある数が選ばれるのだから，その数が指定した $x$ であ

%ここに図

を満たすこととなる。したがって，我々が確率を「測る」方法は，必ずしも，異 なった集合に異なった確率を与えるとは限らないことになる。実際，無限集合 を扱うときには, このことは一般には期待できない.

この議論をもう少し進めて, 実数の有限集合 $A=\left\{x_{1}, x_{2}, \ldots, x_{n}\right\}$ の要素の一 つが選ばれる確率もまた０でなければならない，ということもわかる。確率 $P(A)$ は本来 $\sum_{i=1}^{n} P\left(\left\{x_{n}\right\}\right)$ に等しくあるべきであるうからである。これを敷衍する と，確率関数 $A \mapsto P(A)$ は有限加法性をもつ，つまり， $A_{1}, A_{2}, \ldots, A_{n}$ が互いに 排他的な集合であるならば,

$$
P\left(\bigcup_{i=1}^{n} A_{i}\right)=\sum_{i=1}^{n} P\left(A_{i}\right)
$$

が成り立つと考えられる。この主張は非常にもっともらしい。そして実際, こ れが確率の計算に意味ある基礎を与えるための，本質的な性質の一つであるこ とが後にわかる.

これらよりも少し難しい主張としては，一様分布の下では，任意の可算無限集合 $\left(\right.$ 例えば $\mathbb{Q}$ )の確率は０である，というものがあるが，これはまざに関数 $\mathbf{1}_{\mathbb{Q}}$ の「グラフの下の面積」を解析することによって示されるのである。これを写像 $A \mapsto P(A)$ の $n \rightarrow \infty$ のときの「連続性」の結果であると解釈することもで きる.つまり, $\mathbb{R}$ の部分集合の列 $\left(A_{i}\right)$ が互いに排他的であるなら,

$$
P\left(\bigcup_{i=1}^{\infty} A_{i}\right)=\lim _{n \rightarrow \infty} \sum_{i=1}^{n} P\left(A_{i}\right)=\sum_{i=1}^{\infty} P\left(A_{i}\right)
$$

となる，と考えるのは自然であるからである。

実際，この条件が実数直線 $\mathbb{R}$ 上のルベーグ測度によって満たされることを，第２章でみるだろう。そして，任意の集合上の抽象的な測度の性質を定義する ときにもこれが用いられる。

確率論に含まれる内容は非常に多岐にわたるため，本書では取り上げなかっ た題材も多い。例えば，教科書 [6], [10] で論じられているような確率論への良 い入門となる有限の標本空間とそのエレガントな組合せ論的アイデアについて は，本書では議論しない．我々は本書を通じて，無限の標本空間上の確率現象を 記述するときに,ルベーグ測度が果たす本質的な役割に集中する。これによっ て，他の教科書にみられるような興味深い例や応用は脇においておくことにな るが、その代わりに，密度をもつ確率変数の理論的基礎の首尾一貫した展開を 与えるだろう。


\end{document}
